\section*{Capítulo \ref{ch:g10:reaction-progress-table} --- La tabla de avance}

\begin{solution}{exo:g10:reaction-progress-table:1}
\ce{H2}: $4.0 - 2x$; \ce{O2}: $3.0 - x$; \ce{H2O}: $2x$. El dihidrógeno
se agota en $x = 2.0$, y el dioxígeno, en $x = 3.0$:
$x_{\max} = \qty{2.0}{mol}$, el dihidrógeno es el limitante. \omterm{def:g10:reaction-progress-table:system}{Estado
final}: \qty{0}{mol} de \ce{H2}, \qty{1.0}{mol} de \ce{O2},
\qty{4.0}{mol} de \ce{H2O}.
\end{solution}

\begin{solution}{exo:g10:reaction-progress-table:2}
El carbono se agotaría en $x = 3.0$, y el dioxígeno, en $x = 2.0$:
$x_{\max} = \qty{2.0}{mol}$, el dioxígeno es el limitante. \omterm{def:g10:reaction-progress-table:system}{Estado final}:
\qty{1.0}{mol} de carbono, nada de dioxígeno, \qty{2.0}{mol} de dióxido
de carbono.
\end{solution}

\begin{solution}{exo:g10:reaction-progress-table:3}
$x_{\max} = \qty{0.30}{mol}$ (el azufre es el limitante). \omterm{def:g10:reaction-progress-table:system}{Estado final}:
\qty{0.20}{mol} de hierro, nada de azufre, \qty{0.30}{mol} de sulfuro de
hierro.
\end{solution}

\begin{solution}{exo:g10:reaction-progress-table:4}
El dinitrógeno se agota en $x = 1.0$, y el dihidrógeno, en
$x = 2.4/3 = 0.80$: $x_{\max} = \qty{0.80}{mol}$. \omterm{def:g10:reaction-progress-table:system}{Estado final}:
\qty{0.20}{mol} de \ce{N2}, nada de
\ce{H2}, $2 \times 0.80 = \qty{1.6}{mol}$ \ce{NH3}.
\end{solution}

\begin{solution}{exo:g10:reaction-progress-table:5}
La (b): $4/2 = 2/1$. En la (a), el monóxido de carbono es el limitante;
en la (c), también.
\end{solution}

\begin{solution}{exo:g10:reaction-progress-table:6}
$n(\ce{CH4}) = 32 / 16.0 = \qty{2.0}{mol}$, así que
$x_{\max} = \qty{2.0}{mol}$ para una \omterm{def:g10:reaction-progress-table:stoichiometric}{mezcla estequiométrica}. Dioxígeno:
$2 \times 2.0 = \qty{4.0}{mol}$, es decir,
$4.0 \times 32.0 = \qty{128}{g}$. Agua:
$\qty{4.0}{mol} \times \qty{18.0}{g/mol} = \qty{72}{g}$.
\end{solution}

\begin{solution}{exo:g10:reaction-progress-table:7}
$n(\ce{C3H8}) = 1000 / 44.0 = \qty{22.7}{mol}$; dióxido de carbono:
$3 \times 22.7 = \qty{68.2}{mol}$; $V = 68.2 \times 24.1 =
\qty{1.64e3}{L}$, alrededor de \qty{1.6}{m^3}.
\end{solution}

\begin{solution}{exo:g10:reaction-progress-table:8}
En $x = \qty{1.0}{mol}$: \qty{1.0}{mol} de \ce{CH4}, \qty{1.0}{mol} de
\ce{O2}, \qty{1.0}{mol} de \ce{CO2}, \qty{2.0}{mol} de \ce{H2O}. El
dioxígeno se agota en $x = \qty{1.5}{mol}$: la reacción se detiene ahí,
así que nunca se alcanzan avances mayores.
\end{solution}

\begin{solution}{exo:g10:reaction-progress-table:9}
$n_0(\ce{Zn}) = 1.30 / 65.4 = \qty{0.0199}{mol}$;
$n_0(\ce{Cu^{2+}}) = 0.10 \times 0.1000 = \qty{0.0100}{mol}$. Los \omterm{def:g9:ions:ion}{iones}
cobre son el limitante: $x_{\max} = \qty{0.0100}{mol}$. Cobre formado:
$0.0100 \times 63.5 = \qty{0.635}{g}$; cinc que queda:
$(0.0199 - 0.0100) \times 65.4 = \qty{0.65}{g}$. No quedan \omterm{def:g9:ions:ion}{iones} cobre:
el color azul ha desaparecido.
\end{solution}

\begin{solution}{exo:g10:reaction-progress-table:10}
$n_0(\ce{Mg}) = 0.48 / 24.3 = \qty{0.0198}{mol}$, que se agotaría en
$x = 0.0198$; $n_0(\ce{H+}) = \qty{0.020}{mol}$, que se agotaría en
$x = 0.010$. El ácido es el limitante: $x_{\max} = \qty{0.010}{mol}$,
así que se obtienen \qty{0.010}{mol} de dihidrógeno,
$0.010 \times 24.1 = \qty{0.24}{L}$.
\end{solution}

\begin{solution}{exo:g10:reaction-progress-table:11}
Una \omterm{def:g10:reaction-progress-table:stoichiometric}{mezcla estequiométrica} con \qty{0.30}{mol} de azufre necesita
\qty{0.30}{mol} de hierro; hay \qty{0.20}{mol} más, es decir,
$0.20 / 0.30 \approx \qty{67}{\percent}$ en exceso.
\end{solution}

\begin{solution}{exo:g10:reaction-progress-table:12}
Cantidades iguales: $m(\ce{Fe})/55.8 = m(\ce{S})/32.1$, con
$m(\ce{Fe}) + m(\ce{S}) = \qty{10.0}{g}$. Así,
$m(\ce{Fe}) = 10.0 \times 55.8 / (55.8 + 32.1) = \qty{6.35}{g}$ y
$m(\ce{S}) = \qty{3.65}{g}$.
\end{solution}

\begin{solution}{exo:g10:reaction-progress-table:13}
\begin{enumerate}
  \item $n_0(\ce{CaCO3}) = 10.0 / 100.1 = \qty{0.0999}{mol}$, el único
        \omterm{def:g7:chemical-reactions:reactant}{reactivo}: $x_{\max} = \qty{0.0999}{mol}$; da \qty{0.0999}{mol}
        de \ce{CaO} y de \ce{CO2}, es decir,
        $0.0999 \times 24.1 = \qty{2.41}{L}$ de dióxido de carbono.
  \item $n_0(\ce{CaO}) = \qty{0.0999}{mol}$, $n_0(\ce{H2O}) = 1.8 / 18.0
        = \qty{0.100}{mol}$: la cal viva es (por poco) el limitante; la
        \omterm{def:g6:pure-substances-and-mixtures:mixture}{mezcla} es casi estequiométrica. \ce{Ca(OH)2}:
        $M = 40.1 + 2 \times (16.0 + 1.0) = \qty{74.1}{g/mol}$, masa
        $0.0999 \times 74.1 = \qty{7.40}{g}$.
\end{enumerate}
\end{solution}

\begin{solution}{exo:g10:reaction-progress-table:14}
$n_0(\ce{H+}) = \qty{0.050}{mol}$ en cada matraz; se agotaría en
$x = \qty{0.025}{mol}$. Magnesio: 0.15, 0.30, 0.45, 0.60, 0.75 y
\qty{0.90}{g} dan 0.0062, 0.0123, 0.0185, 0.0247, 0.0309 y
\qty{0.0370}{mol}. En los cuatro primeros matraces el magnesio es el
limitante y el dihidrógeno es $n(\ce{Mg}) \times 24.1$: 0.15, 0.30, 0.45
y \qty{0.60}{L}. En los dos últimos el limitante es el ácido: $0.025 \times 24.1
= \qty{0.60}{L}$. La gráfica es una recta que pasa
por el origen hasta unos \qty{0.61}{g} de magnesio, y después una
meseta horizontal en \qty{0.60}{L}.
\end{solution}

\begin{solution}{exo:g10:reaction-progress-table:15}
$n_0(\text{etanol}) = 4.6 / 46.0 = \qty{0.100}{mol}$ (se agotaría en
$x = 0.100$); $n_0(\ce{O2}) = 4.8 / 24.1 = \qty{0.199}{mol}$ (se
agotaría en $x = 0.199/3 = 0.0664$). El dioxígeno es el limitante,
$x_{\max} = \qty{0.0664}{mol}$. Etanol que queda: $0.100 - 0.0664 =
\qty{0.0336}{mol}$, es decir, $0.0336 \times 46.0 = \qty{1.5}{g}$.
Dióxido de carbono: $2 \times 0.0664 = \qty{0.133}{mol}$, es decir,
$0.133 \times 24.1 = \qty{3.2}{L}$.
\end{solution}

\begin{solution}{pb:g10:reaction-progress-table:1}
\textbf{1.} Izquierda: 2 Na, 6 N. Derecha: 2 Na, $3 \times 2 = 6$ N.
Ajustada.

\textbf{2.} Izquierda: 10 Na, 2 K, 2 N, 6 O. Derecha: K: 2; Na:
$5 \times 2 =
10$; O: $1 + 5 = 6$; N: 2. Ajustada.

\textbf{3.} El capítulo de la \omterm{def:g9:periodic-table-first-look:periodic-table}{tabla periódica}
(\cref{ch:g9:periodic-table-first-look}): el sodio reacciona
violentamente con el agua, dando una disolución corrosiva y
dihidrógeno. Después del choque, el airbag está en contacto con la piel,
los ojos y el \omterm{def:g4:air-a-mixture-of-gases:air}{aire} húmedo de la respiración.

\textbf{4.} $M(\ce{NaN3}) = 23.0 + 3 \times 14.0 = \qty{65.0}{g/mol}$;
$M(\ce{KNO3}) = 39.1 + 14.0 + 3 \times 16.0 = \qty{101.1}{g/mol}$.

\textbf{5.} \ce{NaN3}: $1.00 - 2x$; \ce{Na}: $2x$; \ce{N2}: $3x$.

\textbf{6.} $1.00 - 2x = 0$ en $x_{\max} = \qty{0.500}{mol}$; la azida
de sodio, el único \omterm{def:g7:chemical-reactions:reactant}{reactivo}, es el limitante.

\textbf{7.} \ce{Na}: $2 \times 0.500 = \qty{1.00}{mol}$; \ce{N2}:
$3 \times 0.500 = \qty{1.50}{mol}$.

\textbf{8.} $1.50 \times 24.1 = \qty{36.2}{L}$.

\textbf{9.} \ce{Na}: $1.00 - 10x$; \ce{KNO3}: $n - 2x$; \ce{K2O}: $x$;
\ce{Na2O}: $5x$; \ce{N2}: $x$.

\textbf{10.} $1.00/10 = n/2$, así que $n = \qty{0.200}{mol}$, es decir,
$0.200 \times 101.1 = \qty{20.2}{g}$.

\textbf{11.} $x_{\max} = \qty{0.100}{mol}$: nada de sodio, nada de
nitrato de potasio, \qty{0.100}{mol} de \ce{K2O}, \qty{0.500}{mol} de
\ce{Na2O} y \qty{0.100}{mol} de \ce{N2}.

\textbf{12.} El sodio se agotaría en $x = 0.100$, y el nitrato de
potasio, en $x = 0.150/2 = 0.075$: el nitrato es el limitante,
$x_{\max} =
\qty{0.075}{mol}$, y quedarían en el airbag
$1.00 - 10 \times 0.075 = \qty{0.25}{mol}$ de sodio metálico, listo para
reaccionar con la humedad.

\textbf{13.} $1.50 + 0.100 = \qty{1.60}{mol}$ de nitrógeno por \omterm{def:g10:the-mole:amount}{mol} de
azida de sodio.

\textbf{14.} $n = 60 / 24.1 = \qty{2.49}{mol}$.

\textbf{15.} $2.49 / 1.60 = \qty{1.56}{mol}$ de azida de sodio.

\textbf{16.} $1.56 \times 65.0 = \qty{101}{g}$.

\textbf{17.} $0.200 \times 1.56 = \qty{0.311}{mol}$, es decir,
$0.311 \times 101.1 = \qty{31.5}{g}$ de nitrato de potasio.

\textbf{18.} $2.49 / 1.50 = \qty{1.66}{mol}$, es decir,
$1.66 \times 65.0 = \qty{108}{g}$: la segunda reacción ahorra unos
\qty{7}{g} de azida de sodio.

\textbf{19.} Una \omterm{def:g10:reaction-progress-table:limiting}{reacción total} solo se detiene cuando se agota su
\omterm{def:g10:reaction-progress-table:limiting}{reactivo limitante}; aquí la azida de sodio es el único \omterm{def:g7:chemical-reactions:reactant}{reactivo}, así que
en el \omterm{def:g10:reaction-progress-table:system}{estado final} su cantidad es $1.00 - 2x_{\max} = 0$. No queda nada
tóxico en el airbag.

\textbf{20.} Unos \qty{101}{g} de azida de sodio llenan un airbag de
\qty{60}{L}.
\end{solution}
